\documentclass[11pt,a4paper]{scrartcl}
\usepackage[a4paper, top=20mm, bottom=20mm, left=22mm, right=22mm]{geometry}
\usepackage{lmodern}
\usepackage[T1]{fontenc}
\usepackage[utf8]{inputenc}
\usepackage[ngerman]{babel}
\usepackage{microtype}
\usepackage{parskip}
\usepackage{enumitem}
\usepackage{booktabs}
\usepackage{array}
\usepackage[hidelinks]{hyperref}
\pagestyle{empty}

\begin{document}

\begin{center}
  {\Large\bfseries Zwischenergebnis: OOM-Grenzen bei 256\,MB}\\[4pt]
  {\normalsize Referenz f\"ur zuk\"unftige Benchmark-L\"aufe}\\[2pt]
  {\small Stand: 4.~M\"arz~2026}
\end{center}

\bigskip

\subsection*{Kontext}

100 RocksDB-Durchl\"aufe in Docker mit 256\,MB RAM, kein Swap, 1 CPU.
Workload: \texttt{readwhilewriting} (gleichzeitiges Lesen und Schreiben).
58 erfolgreich, 42 vom OOM-Killer beendet (Exit-Code 137).

\subsection*{Empirisches Speichermodell}

Aus den 100 Durchl\"aufen l\"asst sich folgendes Modell ableiten:

\begin{center}
\fbox{\parbox{0.7\textwidth}{\centering
\textbf{Speicherverbrauch} $\approx$ 2 $\times$ \texttt{write\_buffer\_size} + \texttt{cache\_size} + $\sim$100\,MB Overhead
}}
\end{center}

Der Faktor~2 kommt daher, dass RocksDB bis zu \texttt{max\_write\_buffer\_number}~(=~2)
Memtables gleichzeitig h\"alt. Der Overhead umfasst Block-Index, Filter-Bl\"ocke,
Compaction-Puffer und allgemeinen Prozess-Overhead.

\subsection*{OOM-Matrix bei 256\,MB}

\begin{center}
\begin{tabular}{l *{5}{c}}
\toprule
 & \multicolumn{5}{c}{\texttt{cache\_size}} \\
\cmidrule(lr){2-6}
\texttt{write\_buffer\_size} & 8\,MB & 16\,MB & 32\,MB & 64\,MB & 128\,MB \\
\midrule
4\,MB   & ok & ok & ok & ok & ok \\
16\,MB  & ok & ok & ok & ok & \textbf{OOM} \\
32\,MB  & ok & ok & ok & ok & \textbf{OOM} \\
64\,MB  & \textbf{OOM} & \textbf{OOM} & \textbf{OOM} & \textbf{OOM} & \textbf{OOM} \\
128\,MB & \textbf{OOM} & \textbf{OOM} & \textbf{OOM} & \textbf{OOM} & \textbf{OOM} \\
\bottomrule
\end{tabular}
\end{center}

\textbf{13 sichere Kombinationen} (wbuf $\times$ cache, ohne Kompression):
\begin{itemize}[nosep]
  \item \texttt{wbuf=4\,MB}: alle Cache-Gr\"o\ss en (8--128\,MB) $\rightarrow$ 5 Kombinationen
  \item \texttt{wbuf=16\,MB}: Cache $\leq$ 64\,MB $\rightarrow$ 4 Kombinationen
  \item \texttt{wbuf=32\,MB}: Cache $\leq$ 64\,MB $\rightarrow$ 4 Kombinationen
\end{itemize}

\subsection*{OOM-Grenze bei 256\,MB}

Die empirische Grenze ist:
\[
2 \times \texttt{wbuf} + \texttt{cache} + 100\,\text{MB} \leq 256\,\text{MB}
\quad\Leftrightarrow\quad
2 \times \texttt{wbuf} + \texttt{cache} \leq 156\,\text{MB}
\]

Konfigurationen, die diese Grenze \"uberschreiten, wurden ausnahmslos vom
OOM-Killer beendet. Konfigurationen unterhalb der Grenze liefen alle
erfolgreich durch.

\subsection*{Prognose f\"ur 512\,MB (Raspberry Pi 4)}

Mit dem gleichen Modell:
\[
2 \times \texttt{wbuf} + \texttt{cache} \leq 412\,\text{MB}
\]

\textbf{27 sichere Kombinationen} (14 neu gegen\"uber 256\,MB):

\begin{center}
\begin{tabular}{l *{6}{c}}
\toprule
 & \multicolumn{6}{c}{\texttt{cache\_size}} \\
\cmidrule(lr){2-7}
\texttt{write\_buffer\_size} & 8\,MB & 16\,MB & 32\,MB & 64\,MB & 128\,MB & 256\,MB \\
\midrule
4\,MB   & \textsf{=} & \textsf{=} & \textsf{=} & \textsf{=} & \textsf{=} & \textbf{neu} \\
16\,MB  & \textsf{=} & \textsf{=} & \textsf{=} & \textsf{=} & \textbf{neu} & \textbf{neu} \\
32\,MB  & \textsf{=} & \textsf{=} & \textsf{=} & \textsf{=} & \textbf{neu} & -- \\
64\,MB  & \textbf{neu} & \textbf{neu} & \textbf{neu} & \textbf{neu} & \textbf{neu} & -- \\
128\,MB & \textbf{neu} & \textbf{neu} & \textbf{neu} & \textbf{neu} & \textbf{neu} & -- \\
\bottomrule
\end{tabular}
\end{center}

{\small \textsf{=} = auch bei 256\,MB sicher, \textbf{neu} = nur bei 512\,MB m\"oglich, -- = voraussichtlich OOM}

\subsection*{Konsequenzen f\"ur weitere Benchmarks}

\begin{enumerate}[nosep]
  \item \textbf{256\,MB Extended Run:} Nur die 13 sicheren Kombinationen mit
        Wiederholungen (3~Reps, 2~Workloads). Erwartete Dauer: $\sim$20\,h mit
        4 parallelen Containern.
  \item \textbf{512\,MB Extended Run:} Nur die 14 \emph{neuen} Kombinationen,
        die bei 256\,MB nicht m\"oglich waren (56 Param-Kombinationen mit Kompression).
        Erwartete Dauer: $\sim$21\,h mit 4 parallelen Containern.
  \item \textbf{OOM-Konfigurationen \"uberspringen:} Konfigurationen, die die
        Speichergrenze \"uberschreiten, werden nicht mehr getestet. Das spart
        $\sim$30\,min pro OOM-Durchlauf (42 OOMs $\times$ 30\,min = 21\,h eingespart).
\end{enumerate}

\subsection*{Bedeutung f\"ur die Bachelorarbeit}

\begin{itemize}[nosep]
  \item Das Speichermodell zeigt, dass \texttt{write\_buffer\_size} der
        \emph{dominante} Speicherverbraucher ist (Faktor~2 durch doppelte Memtables).
  \item F\"ur NebulaStream auf 256\,MB-Ger\"aten darf der Schreibpuffer maximal
        32\,MB gro\ss{} sein -- idealerweise 4\,MB f\"ur maximale Stabilit\"at.
  \item Die Cache-Gr\"o\ss e kann flexibler gew\"ahlt werden, muss aber zusammen
        mit dem Schreibpuffer die Gesamtgrenze einhalten.
  \item Auf moderneren Edge-Ger\"aten (512\,MB, z.\,B.\ Raspberry Pi 4) sind
        gr\"o\ss ere Schreibpuffer (64--128\,MB) m\"oglich, was die
        Schreibverst\"arkung deutlich reduzieren sollte.
\end{itemize}

\end{document}
